\section{System and Threat Model}
\label{sec:system}

\subsection{Personalization lifecycle}

The device runs a frozen or infrequently updated base SLM and maintains two mutable forms of personalization. First, a local memory service stores selected records and retrieves them into inference context. Second, a trainer periodically forms a window $W_k$ and produces a candidate LoRA adapter $q_k$. A promotion pointer selects at most one candidate as the active adapter. The two sinks may operate at different cadences, but they share a lineage service and policy state.

\Cref{fig:architecture} separates the data plane from the control plane. Raw content is captured at a trusted ingress boundary. The provenance tagger canonicalizes the payload, retrieves OS-observable attributes, and creates an authenticated envelope. The policy engine issues sink-specific decisions. Derived summaries inherit all parents rather than receiving a fresh ``trusted'' label. The memory service and trainer emit lineage edges transactionally with their state changes. The audit harness evaluates a candidate without exposing it as the active model. Promotion commits the candidate, audit report, and lineage root atomically.

\subsection{Trust assumptions}

\begin{assumption}[Trusted computing base]
\label{ass:tcb}
The mobile OS kernel, provenance service, secure storage used for graph roots, base-model loader, policy engine, lineage recorder, and training runtime execute as specified. Untrusted applications cannot forge another app identity or rewrite a committed event envelope.
\end{assumption}

\begin{assumption}[Complete mediation]
\label{ass:mediation}
Every state-changing operation on memory, training data, candidate adapters, audit decisions, and promoted checkpoints is mediated by the policy/lineage layer. There is no side channel that writes persistent learned state without producing the corresponding lineage edges.
\end{assumption}

\begin{assumption}[Append-only event identity]
\label{ass:identity}
Each event and derived object has a collision-resistant content identifier over its canonical payload, provenance envelope, parent identifiers, and version. Mutation creates a new object rather than changing the old identifier.
\end{assumption}

These assumptions are strong but inspectable. If a vendor allows an application to write directly to the training cache, or if a summarizer discards its parent list, the main guarantee does not apply. PoisonGuard is therefore a system architecture, not a classifier that compensates for arbitrary implementation bypasses.

\subsection{Adversary}

The adversary controls content delivered through one or more non-user channels. It may send messages or calendar invitations, trigger notifications through a malicious app, share files, place text on a synchronized clipboard, or induce content that later appears in a screenshot index. It knows the broad policy and may adapt payloads after observing normal model behavior. It can attempt a memory attack, a training-data backdoor, or a two-stage laundering attack.

The adversary may also socially engineer the user into reproducing a trigger in typed text. This capability is treated separately because the resulting OS evidence legitimately records user authorship. A defense that silently equates ``typed'' with ``benign'' would conceal this residual risk. Our evaluation therefore reports direct-channel and induced-user attacks separately.

The attacker cannot compromise the TCB in \cref{ass:tcb}, forge event authentication, extract secure storage, alter the experiment harness, or physically modify the device. Denial of service through unbounded incoming content is rate-limited but is not the central goal.

\subsection{Assets and security properties}

The principal asset is behavioral integrity of the personalized SLM. Confidentiality and privacy remain important but are not the theorem target. We require four properties.

\begin{definition}[Provenance preservation]
For every persistent object $v$, the graph contains a path from each direct or indirect source event that influenced $v$ to $v$.
\end{definition}

\begin{definition}[Non-laundering]
If policy denies event $e$ for adapter training, then no promoted checkpoint may be a descendant of $e$. If $e$ is later revoked, no checkpoint descendant may remain active after recovery completes.
\end{definition}

\begin{definition}[Quarantine safety]
A candidate adapter cannot serve ordinary user inference before a successful audit and an atomic promotion that records its lineage closure.
\end{definition}

\begin{definition}[Counterfactual restoration]
After revoking $R_t$, the recovered state equals the state obtained by executing the same versioned personalization pipeline on the admissible event subsequence with the revoked lineage closure removed.
\end{definition}

Counterfactual restoration is stronger than ``attack success decreased'' and narrower than semantic repair. Equality may be bitwise for deterministic integer kernels or within a declared tolerance for deterministic floating-point schedules. The experiment must state which form it checks.

\subsection{Deployment boundary}

The design fits a system service, a privileged personalization framework, or a research runtime. A third-party application can implement a weaker variant, but cannot reliably attribute cross-application content without OS support. For that reason, application identity and user-authorship evidence are observations from the trusted ingress boundary rather than strings supplied by the text producer.

\paragraph{Topic-alignment checkpoint.} The architecture remains single-device and mobile. Federated workers, cloud aggregation, and differential privacy are outside the state machine.

\begin{figure}[t]
  \centering
  \resizebox{0.95\columnwidth}{!}{\begin{tikzpicture}[
  hnode/.style={honest, font=\small, text width=32mm, minimum height=9.5mm, inner sep=3pt},
  anode/.style={attack, font=\small, text width=32mm, minimum height=9.5mm, inner sep=3pt},
  dnode/.style={defense, font=\small, text width=72mm, minimum height=9.5mm, inner sep=3pt},
  inode/.style={infra, font=\small, text width=32mm, minimum height=9.5mm, inner sep=3pt}
]
  % Tier 1: Ingress Sources (Side-by-side, y = 0)
  \node[hnode] (user) at (0, 0) {\tikz\draw[pgblue,line width=.9pt] (0,.15) circle (.12) (-.2,-.25) arc (180:0:.2 and .18);\\\textbf{User-Typed Input}\\[0.3mm]\scriptsize Direct human entry};
  \node[anode] (attacker) at (4.0, 0) {\tikz\draw[pgorange,line width=.9pt] (0,.15) circle (.12) (-.2,-.25)--(0,-.01)--(.2,-.25);\\\textbf{Remote / App Attacker}\\[0.3mm]\scriptsize Untrusted channels};

  % Tier 2: Policy & Mediation Gate (y = -2.0) -> Arrow length = 11mm!
  \node[dnode] (guard) at (2.0, -2.0) {\textbf{Origin Binding \& Sink-Specific Policy Gate}\\[0.3mm]\scriptsize Complete OS mediation: binds cryptographic provenance to all events};

  % Tier 3: Persistent Sinks (Side-by-side, y = -3.8) -> Arrow length = 9mm!
  \node[hnode] (mem) at (0, -3.8) {\textbf{Agent Memory Store}\\[0.3mm]\scriptsize $\pi_{\text{mem}}$: allow / quarantine};
  \node[inode] (lora) at (4.0, -3.8) {\textbf{LoRA Trainer}\\[0.3mm]\scriptsize $\pi_{\text{train}}$: quarantine / deny};

  % Tier 4: Unified Defense Shield (y = -5.4) -> Arrow length = 9mm!
  \node[dnode] (lineage) at (2.0, -5.4) {\textbf{Shared Lineage \& Rollback Shield}\\[0.3mm]\scriptsize Causal graph tracking, adapter quarantine, and deterministic rollback};

  % Prominent, long arrows with clearly offset labels (NO OVERLAPPING!)
  \draw[attackflow] (attacker.north) to[bend right=25] node[above=1.5mm,font=\scriptsize]{induced action (phishing)} (user.north);

  % Ingress into Policy Gate: 11mm long vertical arrows!
  \draw[dataflow] (user.south) -- (user.south |- guard.north) node[midway,left=2mm,font=\small\bfseries,text=pgblue]{trusted};
  \draw[attackflow] (attacker.south) -- (attacker.south |- guard.north) node[midway,right=2mm,font=\small\bfseries,text=pgorange]{untrusted};

  % Policy Gate into Sinks: 9mm long vertical arrows!
  \draw[dataflow] (mem.north |- guard.south) -- (mem.north) node[midway,left=2mm,font=\small\bfseries,text=pgblue]{$\pi_{\text{mem}}$};
  \draw[dataflow] (lora.north |- guard.south) -- (lora.north) node[midway,right=2mm,font=\small\bfseries,text=pggray]{$\pi_{\text{train}}$};

  % Sinks into Lineage Shield: 9mm long vertical arrows!
  \draw[controlflow] (mem.south) -- (mem.south |- lineage.north) node[midway,left=2mm,font=\small\bfseries,text=pggreen]{audit};
  \draw[controlflow] (lora.south) -- (lora.south |- lineage.north) node[midway,right=2mm,font=\small\bfseries,text=pggreen]{commit};
\end{tikzpicture}
}
  \caption{Threat and control boundary. Orange patterned sources are untrusted; blue sources are directly user-authored; green shield stages mediate both persistent sinks. The adaptive retyping path is intentionally shown outside the automatic provenance guarantee.}
  \label{fig:threat}
\end{figure}
